\section{Preliminaries}
\label{sec:preliminaries}

This section defines notation and gives interface-level definitions for the lattice primitives, the NIZK model, the
blind-signature syntax, ARC syntax and security goals, and pseudorandom functions.



\subsection{Notation}
Let $\lambda$ denote the security parameter and  \(\negl(\lambda)\) be any positive function negligible in  \(\lambda\). We write PPT for
probabilistic algorithms whose runtime is polynomial in \(\lambda\). For integers \(a<b\) we write
\([a,b]\coloneqq\{a,a+1,\ldots,b\}\) and we use \([b]\) as shorthand for \([1,b]\). Throughout this work, we denote vectors in lower-case bold letters such as $\vectorify{b}$ and matrices in upper-case bold letters such as $\mA$.

We work over the cyclotomic ring \(\R=\Z[X]/(X^N+1)\) for power-of-two \(N\) and its quotient \(\Rq=\R/q\R\), identified
with polynomials of degree \(<N\) with coefficients in \(\Z_q\). For \(a\in\Z\), we write \(\langle a\rangle_q\in\{0,\dots,q-1\}\)
for reduction modulo \(q\). When considering $\Z_q$-elements from an interval \([-a,b]\), we identify these with the set  $\{q-a,q-a+1,\dots,0,1,\dots,b\}$.
For elements $x\in R_q$ we write $||x||\leq B$ to mean that each coefficient of $x$ comes from $[-B,B]$. For vectors $\vectorify{x}\in R_q^\ell$, $||\vectorify{x}||\leq B$ means that the statement holds for all elements of $\vectorify{x}$.

\subsection{Commitment schemes and Ajtai commitments}
\label{subsec:prelim-commit}
\label{subsec:commitments}

\begin{definition}[Commitment scheme]\label{def:commitment-scheme}
Let $\mathcal{M}$ be a finite message space. A (non-interactive) commitment scheme is a tuple of PPT algorithms
\[
  \ComScheme=(\ComSetup,\Commit,\ComVerify)
\]
with the following syntax:
\begin{itemize}[leftmargin=1.25em]
  \item $\ComPublicParam\gets \ComSetup(\SecParam)$ is a PPT algorithm which outputs public parameters $\ComPublicParam$. It is implicit input to all other algorithms.
  \item $c \gets \Commit(m)$ is a PPT algorithm which takes a message $m\in\mathcal{M}$ as input and outputs a commitment $c$ and decommitment $d$.
  %\item $\pi \gets \Open(\ComPublicParam,m;r)$ outputs an opening string (often just $r$). \cb{TODO: remove open}
  \item $\ComVerify(c,m,d)\in\{0,1\}$ deterministically checks validity of the opening.
\end{itemize}
Correctness holds if for all $\ComPublicParam\gets \ComSetup$,  $m\in \mathcal{M}$ and $(c,d)\gets \Commit(m)$ it holds that
\[
  \ComVerify(m,c,d)=1.
\]

The scheme is \emph{computationally binding} if for all PPT adversaries $\Adv$,
\[
  \Pr_r\Big[\begin{matrix}
  	(\ComPublicParam,c,m,d,m',d')\gets \Adv(r) \wedge m\neq m' \\
  	 \ComVerify(c,m,d)=1 \wedge \ComVerify(c,m',d')=1 
  \end{matrix} 
   \Big]\le \negl(\lambda)\text{.}
\]

The scheme is \emph{(computationally) hiding} if for all PPT distinguishers $\mathcal{D}$, all 
messages $m_0,m_1\in\mathcal{M}$ and $\ComPublicParam\gets \ComSetup()$,
\[
  \Big|\Pr[\mathcal{D}^{\mathcal{O}_\Commit(m_0)}(\ComPublicParam)=1]
  -\Pr[\mathcal{D}^{\mathcal{O}_\Commit(m_1)}(\ComPublicParam)=1]\Big|\le \negl(\lambda)
\]
where $\mathcal{O}_{\Commit}(m_b)$ runs $(c,d)\gets \Commit(m_b)$ and outputs $c$ and the probability is taken over the randomness of $\Commit,\mathcal{D}$.
\end{definition}

\subsubsection{Ajtai's Commitment scheme}
We use the standard bounded linear commitment over \(\Rq\), matching the compact Ajtai-style commitment of
\cite[Definition~13]{EPRINT:AttCraKoh21}. Fix message length \(\ell_m\), randomness dimension \(\ell_r\), output length
\(\ell_{\ComScheme}\), and coefficient bound \(\BoundMsg\), and let \(\ell_{\mathsf{in}}=\ell_m+\ell_r\).
\begin{itemize}
  \item \(\ComSetup(\SecParam)\) samples \(\mACom\getsunif \Rq^{\ell_{\ComScheme}\times \ell_{\mathsf{in}}}\) and outputs
  \(\ComPublicParam:=(\mACom,\BoundMsg)\).
  \item The message space is
  \[
    \mathcal{M}:=\{\vectorify{m}\in\Rq^{\ell_m}\mid \|\vectorify{m}\|_\infty\le \BoundMsg\}.
  \]
  \item \(\Commit(\vectorify{m})\) samples
  \(\vectorify{r}\getsunif[-\BoundMsg,\BoundMsg]^{\ell_r}\subseteq\Rq^{\ell_r}\) and outputs
  \[
    \vectorify{c}=\mACom\cdot[\vectorify{m}\Vert \vectorify{r}]^\top
  \]
  together with opening \(\vectorify{r}\).
  \item \(\ComVerify(\vectorify{c},\vectorify{m},\vectorify{r})\) accepts iff
  \[
    \vectorify{c}\in\Rq^{\ell_{\ComScheme}},\qquad
    \vectorify{m}\in\mathcal{M},\qquad
    \|\vectorify{r}\|_\infty\le \BoundMsg,
  \]
  and \(\vectorify{c}=\mACom\cdot[\vectorify{m}\Vert \vectorify{r}]^\top\).
\end{itemize}

\begin{definition}[Module-SIS and module-ISIS over \(R_q\)]\label{def:sis-isis}
Fix integers \(n,m\ge 1\), modulus \(q\ge 2\), and bound \(\beta>0\).
\begin{itemize}[leftmargin=1.25em]
  \item \textbf{\(\mathsf{MSIS}_\infty(n,m,q,\beta)\).} Given \(\mA\in R_q^{n\times m}\), find a non-zero
  \(\vectorify{s}\in R^m\) such that \(\mA\vectorify{s}\equiv 0 \pmod q\) and \(\|\vectorify{s}\|_\infty\le \beta\).
  \item \textbf{\(\mathsf{MISIS}_\infty(n,m,q,\beta)\).} Given \((\mA,\vect)\in R_q^{n\times m}\times R_q^n\), find
  \(\vectorify{s}\in R^m\) such that \(\mA\vectorify{s}\equiv \vect \pmod q\) and
  \(\|\vectorify{s}\|_\infty\le \beta\).
\end{itemize}
\end{definition}

\begin{lemma}[Binding of the Ajtai commitment]\label{lem:ajtai-binding}
Suppose an adversary outputs a commitment \(\vectorify{c}\) and two distinct valid openings
\((\vectorify{m},\vectorify{r})\neq(\vectorify{m}',\vectorify{r}')\) with coefficient bound \(\BoundMsg\). Then
\[
  \vectorify{\Delta}:=[\vectorify{m}-\vectorify{m}'\Vert \vectorify{r}-\vectorify{r}']^\top
\]
is a non-zero solution to \(\mathsf{MSIS}_\infty\) for \(\mACom\) with bound \(2\BoundMsg\), since
\(\mACom\vectorify{\Delta}\equiv 0\pmod q\) and \(\|\vectorify{\Delta}\|_\infty\le 2\BoundMsg\). Hence the commitment is
computationally binding whenever \(\mathsf{MSIS}_\infty\) is hard for the resulting parameters
\cite[Lemma~2]{EPRINT:AttCraKoh21}.
\end{lemma}

\begin{lemma}[Statistical hiding of the Ajtai commitment]\label{lem:ajtai-hiding}
Write \(\mACom=[\mA_{\mathsf{msg}}\mid \mA_{\mathsf{rnd}}]\) according to the message and randomness blocks. If the
randomness dimension \(\ell_r\) satisfies the leftover-hash entropy condition of
\cite[Lemma~1]{EPRINT:AttCraKoh21} for coefficient alphabet \([-\BoundMsg,\BoundMsg]\), then
\(\mA_{\mathsf{rnd}}\vectorify{r}\) is statistically close to uniform over \(\Rq^{\ell_{\ComScheme}}\). Consequently,
\[
  \mACom\cdot[\vectorify{m}\Vert \vectorify{r}]^\top
  = \mA_{\mathsf{msg}}\vectorify{m}+\mA_{\mathsf{rnd}}\vectorify{r}
\]
statistically hides \(\vectorify{m}\).
\end{lemma}

\subsection{Lattices, trapdoors, vSIS, and rational hashing}
\label{subsec:prelim-vsis}
We recall only the definitions needed to state our construction; implementation details of the trapdoor sampler are
deferred to Appendix~\ref{app:gaussian-sampler}.

\begin{definition}[Lattice trapdoor suite]\label{def:trapdoor-suite}
A lattice trapdoor suite over \(\Rq\) consists of PPT algorithms
\[
  (\TrapGen,\SampPre)
\]
and admissible public parameters \((n,m,q,\BoundSig)\) as in \cite[Definition~1]{EPRINT:DKLW25} such that:
\begin{itemize}[leftmargin=1.25em]
  \item \((\mA,\trapdoor)\gets \TrapGen(\SecParam)\) outputs a public matrix \(\mA\in\Rq^{n\times m}\) and a trapdoor
  \(\trapdoor\), where the distribution of \(\mA\) is statistically close to uniform over \(\Rq^{n\times m}\).
  \item On input \((\trapdoor,\vect)\) with \(\vect\in\Rq^n\), the algorithm \(\SampPre(\trapdoor,\vect)\) outputs
  \(\vecu\in R^m\) such that \(\mA\vecu=\vect\) in \(\Rq^n\), except with negligible failure probability, and
  \(\|\vecu\|_\infty\le \BoundSig\) except with negligible probability
  \cite{STOC:GenPeiVai08,EPRINT:DKLW25}.
\end{itemize}
\end{definition}

\begin{definition}[Trapdoor rational hash]\label{def:bbs-rational-hash}
A trapdoor rational hash scheme over \(\Rq\) is a tuple
\[
  \TDRH=(\Setup,\TrapGen,\mathsf{Eval},\SampPre)
\]
consisting of the lattice trapdoor suite above together with a public evaluation algorithm. On input \(\SecParam\),
\(\Setup\) outputs a public hash description \(B\), a message space \(\mathcal{M}\), a randomness space \(\mathcal{X}\), an
admissibility error bound \(\delta\), and public bounds \((\BoundMsg,\BoundSig)\). For fixed \(B\), the algorithm
\(\mathsf{Eval}(B,m,\chi)\) takes a message \(m\in\mathcal{M}\) and witness randomness \(\chi\in\mathcal{X}\) and outputs a
target \(\vect\in\Rq^n\) or \(\bot\) if the input is inadmissible.

For a public key \(B\), define the admissible message subspace
\[
  \mathcal{M}_B:=\left\{m\in\mathcal{M} :
  \Pr_{\chi\gets \mathcal{X}}\big[\mathsf{Eval}(B,m,\chi)=\bot\big]\le \delta\right\},
\]
following the \(\Sigma_{\mathsf{vSIS}}\) template of \cite[Section~4.1]{EPRINT:DKLW25}. Correctness requires that for
every \(m\in\mathcal{M}_B\) and every \(\chi\in\mathcal{X}\) with \(\mathsf{Eval}(B,m,\chi)=\vect\neq\bot\), the output
\(\vecu\gets \SampPre(\trapdoor,\vect)\) satisfies \(\mA\vecu=\vect\) and \(\|\vecu\|_\infty\le \BoundSig\), except with
negligible probability.
\end{definition}

\begin{definition}[ARC rational hash and cleared-denominator form]\label{def:rational-hash}
Our concrete instantiation uses a public key \(B=(B_0,B_1,B_2,B_3)\in\Rq^4\), message \(m\in\Rq\), and witness
randomness \(\chi=(r_0,r_1)\in\Rq^2\) with coefficient bounds
\(\|m\|_\infty,\|r_0\|_\infty,\|r_1\|_\infty\le \BoundMsg\). The evaluation algorithm is
\[
  h_{m,(r_0,r_1)}(B):=
  \begin{cases}
    \bigl(B_0+B_1\cdot m+B_2\cdot r_0\bigr)\ \Hdiv\ (B_3-r_1)\\
    \qquad\text{if \(B_3-r_1\) is invertible in every slot,}\\
    \bot \qquad\text{otherwise.}
  \end{cases}
\]
Whenever \(\vect=h_{m,(r_0,r_1)}(B)\) is defined, it is equivalently characterized by the cleared-denominator identity
\begin{equation}\label{eq:hash-cleared}
  (B_3-r_1)\Had \vect-\bigl(B_0+B_1\cdot m+B_2\cdot r_0\bigr)=0.
\end{equation}
In the implementation, the same abstract message \(m\) is stored in packed rows \((M,K)\) with \(m=M+K\), so
the same identity is written as
\[
  (B_3-R_1)\Had \vect-\bigl(B_0+B_1\cdot(M+K)+B_2\cdot R_0\bigr)=0.
\]
\end{definition}

\begin{remark}[Relation to the DKLW25 BBS candidate]\label{rem:hash-notation-match}
Our ARC family is exactly the BBS-type rational hash studied in \cite[Table~1]{EPRINT:DKLW25}, specialized by
\[
\begin{aligned}
  \tilde b_0&=(B_0,B_1,B_2),\qquad \tilde b_1=B_3,\\
  u&=m,\qquad x_0=r_0,\qquad x_1=r_1.
\end{aligned}
\]
Equivalently,
\[
  h_{m,(r_0,r_1)}(B)
  =\frac{B_0+B_1m+B_2r_0}{B_3-r_1}.
\]
The undefined case \(h_{m,(r_0,r_1)}(B)=\bot\) is exactly the inadmissibility event that \(B_3-r_1\) is not invertible
slotwise, and is absorbed by the admissibility parameter \(\delta\). In issuance, \((r_0,r_1)\) is obtained from the
holder and issuer contributions via the centering map of Appendix~\ref{app:defs}; by
Lemma~\ref{lem:center-uniform}, this is compatible with the bounded randomness space used in the DKLW security
experiments.
\end{remark}

\begin{definition}[vSIS-style hash-and-preimage signature template]\label{def:vsis-template}
Fix a trapdoor rational hash scheme \(\TDRH\) with public output
\((B,\mathcal{M},\mathcal{X},\delta,\BoundMsg,\BoundSig)\). The associated signature scheme
\(\vsisscheme=(\KeyGen,\Sign,\SigVerify)\) is the \(\Sigma_{\mathsf{vSIS}}\) template of
\cite[Figure~4 and Section~4.1]{EPRINT:DKLW25}, specialized to the public function description \(B\) and the
evaluation map \(\mathsf{Eval}\):
\begin{itemize}[leftmargin=1.25em]
  \item \(\KeyGen(\SecParam)\) runs \((\mA,\trapdoor)\gets \TrapGen(\SecParam)\) and outputs
  \(\vk=\mA\), \(\sk=\trapdoor\).
  \item \(\Sign(\sk,m)\) for \(m\in\mathcal{M}_B\) samples \(\chi\gets\mathcal{X}\) until
  \(\vect=\mathsf{Eval}(B,m,\chi)\neq\bot\), computes \(\vecu\gets \SampPre(\trapdoor,\vect)\), and outputs
  \(\sigma=(\vecu,\chi)\).
  \item \(\SigVerify(\vk,m,\sigma)\) for \(\sigma=(\vecu,\chi)\) accepts iff \(m\in\mathcal{M}_B\),
  \(\chi\in\mathcal{X}\), \(\mathsf{Eval}(B,m,\chi)\neq\bot\), \(\|\vecu\|_\infty\le \BoundSig\), and
  \[
    \mA\vecu=\mathsf{Eval}(B,m,\chi).
  \]
\end{itemize}
\end{definition}

\begin{definition}[Hinted-vSIS and \(\mathsf{GenISIS}_f\) variants from DKLW25]\label{def:dklw-assumptions}
We recall the DKLW25 assumption variants used in the reduction discussion below.
Fix a public left factor \(F\), hint space \(H\), target space \(G\), predicate \(P\), and query bound \(Q\), all as in
\cite[Definition~5]{EPRINT:DKLW25}.
\begin{itemize}[leftmargin=1.25em]
  \item \textbf{\(\mathsf{Hint\mbox{-}vSIS}\).} The adversary chooses a query set \(\mathcal{Q}\subseteq H\) and a fresh
  target function \(g^\star\in G\setminus \mathcal{Q}\) before seeing the sampled matrix \(A\). It then receives short
  preimages of the queried hints \(q(A)\) under \(F(A)\) for every \(q\in\mathcal{Q}\), and must output a non-zero short
  vector \(u^\star\) such that \(F(A)u^\star=g^\star(A)\). The instance is restricted by the predicate condition
  \(P(F\cup \mathcal{Q}\cup(\{g^\star\}\setminus\{0\}))=1\).
  \item \textbf{\(s\)-\(\mathsf{Hint\mbox{-}vSIS}\).} This is the strong variant of the same experiment. The adversary still
  chooses the hint set \(\mathcal{Q}\) before seeing \(A\), but it is allowed to choose the fresh target \(g^\star\) only
  after seeing \(A\) and the hinted preimages.
  \item \textbf{\(s\)-\(\$\)-\(\mathsf{Hint\mbox{-}vSIS}\).} This is the strong random-hinted variant. It is the same as
  \(s\)-\(\mathsf{Hint\mbox{-}vSIS}\), except that the query set \(\mathcal{Q}\) is sampled uniformly at random from \(H\)
  rather than chosen by the adversary \cite[Definition~5]{EPRINT:DKLW25}.
\end{itemize}

Now fix a keyed evaluation family \(f:K\times M\times X\to R_q^n\) as in \cite[Definition~8]{EPRINT:DKLW25}.
\begin{itemize}[leftmargin=1.25em]
  \item \textbf{\(\mathsf{GenISIS}_f\).} The challenger samples a public key \(\kappa\gets K\) and a uniform matrix
  \(A\getsunif R_q^{n\times m}\). It provides a preimage oracle that, on input \(\mu\in M\), samples \(\chi\gets X\),
  sets \(t=f(\kappa,\mu,\chi)\), returns a short preimage \(s\) such that \(As=t\), and records the tuple
  \((s,\mu,\chi)\). The adversary wins if it outputs a fresh tuple \((s^\star,\mu^\star,\chi^\star)\) not previously
  returned by the oracle such that \(As^\star=f(\kappa,\mu^\star,\chi^\star)\) and \(0<\|s^\star\|\le \beta\).
  \item \textbf{\(\mathsf{IntGenISIS}_f\).} DKLW25 also defines an interactive extension in which the challenger offers a
  signing-style oracle with an additional linear masking term, and Theorem~6 shows how this interactive problem reduces
  back to \(\mathsf{GenISIS}_f\) under enlarged parameters.
\end{itemize}
\end{definition}

\begin{remark}[Security interpretation via DKLW25]\label{rem:vsis-unforgeability}
For our family, the DKLW25 BBS row is instantiated by
\[
  f(B,m,\chi)=\mathsf{Eval}(B,m,\chi)=h_{m,(r_0,r_1)}(B)
\]
with \(u=m\), \(x_0=r_0\), \(x_1=r_1\), \(\tilde b_0=(B_0,B_1,B_2)\), and \(\tilde b_1=B_3\). Assume negligible
inadmissibility \(\delta\), coefficient-wise bounds on the denominator data, and the natural-predicate condition of
\cite[Definition~7]{EPRINT:DKLW25} under the same split-ring and entropy hypotheses used in DKLW's BBS discussion.
Under these conditions, Theorem~4 gives correctness of the non-blind \(\Sigma_{\mathsf{vSIS}}\) template instantiated by
our family. Theorem~5 then yields \(s\)EUF-RMA of that non-blind signature layer from the strong random-hinted vSIS
assumption once the predicate condition holds. Theorem~7 shows that the same non-blind security experiment is equivalent
to \(\mathsf{GenISIS}_f\) for \(f(B,m,\chi)=\mathsf{Eval}(B,m,\chi)\), up to the same trapdoor-to-uniform hybrids made
explicit in Appendix~A.3 of \cite{EPRINT:DKLW25}. Theorem~8 and Corollary~1 then relate \(\mathsf{GenISIS}_f\) and the
strong random-hinted vSIS assumption for this family, with Appendix~A.4 giving the converse direction. This is
therefore the correct assumption-level justification for the underlying non-blind rational-hash signature family used
here. It does not prove message hiding of \(h_{m,\chi}(B)\). It also does not yet prove one-more unforgeability or
blindness of the blind issuance protocol of Section~\ref{sec:blind-signature}.
\end{remark}

\subsection{Non-interactive zero-knowledge arguments of knowledge (NIZKs)}
\label{subsec:prelim-nizk}
\label{subsec:nizk-ark}
We use the NIZK-ARK syntax of the SmallWood framework instantiated in the random-oracle model
\cite[Definition~3]{EPRINT:FenRiv25b}.

\begin{definition}[NIZKs in the random oracle model]\label{def:nizk-ark}
Let $\mathcal{R}\subseteq \mathcal{X}\times \mathcal{W}$  be a family of relations.
A (transparent) non-interactive zero-knowledge \emph{argument of knowledge} for $\mathcal{R}$ in the random oracle model
(with oracle $\RO(\cdot)$) is a pair of algorithms $(\ZKProve^{\RO},\ZKVerify^{\RO})$:
\begin{itemize}[leftmargin=1.25em]
  \item $\ZKProve^{\RO}(w ,x)\to \pi$ is probabilistic and on input $(x,w)$ with
  $(x,w)\in R$ outputs a proof $\pi$.
  \item $\ZKVerify^{\RO}(\pi,x)\to \{0,1\}$ is a deterministic polynomial-time algorithm.
\end{itemize}

\paragraph{Completeness.}
For all $(x,w)\in\mathcal{R}$,
\[
  \Pr\big[\ZKVerify^{\RO}(\pi,x)=1 \ :\ \pi\gets \ZKProve^{\RO}(w,x)\big]=1.
\]

\paragraph{Straightline-extractable knowledge soundness.} 
There exists a PPT algorithm $\Ext$, also called \emph{extractor}, such that for any PPT adversary $\Adv$,
\[
  \Pr\Big[
    \ZKVerify^{\RO}(\pi,x)=1 \ \wedge\ (x,w)\notin R
  \Big]\ \le\ \varepsilon,
\]
where $(\pi,Q)\leftarrow \Adv^{\RO}(x)$, $Q$ is the set of query--response pairs made to $\RO$ by $\Adv$,
and $w\leftarrow \Ext(Q,\pi)$.

\paragraph{Zero knowledge.}
There exists a PPT algorithm $\Sim$, called the \emph{simulator}, such that for any $(x,w)\in \mathcal{R}$ and PPT algorithm $\Adv$:
\[
  \Big|\Pr[\Adv^{\RO}(\pi_{\mathsf{real}})=1]-\Pr[\Adv^{\RO}(\pi_{\mathsf{sim}})=1]\Big|
  \le \xi,
\]
where $\pi_{\mathsf{real}}\gets \ZKProve^{\RO}(w ,x)$,  $\pi_{\mathsf{sim}}\gets \Sim^{\RO}(x)$ and $\Sim$ can program the output of $\RO$ on any input.
\end{definition}

We instantiate these proofs using the SmallWood PACS/PIOP framework compiled with Fiat--Shamir in the random oracle
model \cite{EPRINT:FenRiv25b}. SmallWood is a commit-and-prove proof system that only relies on symmetric primitives. While its proof size is linear in $|w|$ it is sublinear in the circuit size $x$.

%At a high level, SmallWood expresses statements as polynomial constraints evaluated
%on an $\Omega\subset\Fq$ packing domain, supports both \emph{parallel} constraints checked point-wise on $\Omega$ and
%\emph{aggregated} constraints checked via an $\Omega$-sum identity, and uses random verifier queries outside $\Omega$ to
%spot-check correctness of committed witness polynomials while preserving HVZK via blinding evaluations
%(Section~\ref{sec:smallwood-model}). 



\subsection{Blind signatures}
A blind signature algorithm allows a holder $\Holder$ to obtain signatures valid under the verification key of an issuer $\Issuer$, without the latter learning the signed message $\mu$. Our definitions follow \cite{C:Fischlin06}. When any algorithm $\mathcal{B}$ uses the interactive protocol $\Issue$ with $\Issuer$, we will write $\mathcal{B}^{\Issue:\langle \Issuer(\cdot):\cdot \rangle}$ and adjust this the notation for other interactions accordingly.

\begin{definition}[Blind Signature]\label{def:bs-syntax}
	A Blind Signature scheme is a tuple of algorithms
	\[
	\BlindSignature=(\Setup,\IKeyGen,\Issue,\Verify)
	\]
	such that:
	\begin{itemize}[leftmargin=1.25em]
		\item $\Setup(\SecParam)$ is a PPT algorithm that outputs public parameters $\PublicParamsBS$. All other algorithms obtain $\SecParam,\PublicParamsBS$ as implicit inputs.
		\item $\IKeyGen()$ is a PPT algorithm that outputs issuer keys $(\vk,\sk)$.
		\item $\Issue$ is an interactive protocol between $\Issuer$ and $\Holder$, where $\Issuer$ has input $\sk$ and $\Holder$ has input $\mu,\vk$. At the end of the protocol, $\Holder$ obtains $\sig$.
		\item $\Verify(\vk,\mu,\sig)$ is a deterministic polynomial-time algorithm that outputs a bit.
	\end{itemize}
\end{definition}

\begin{definition}[Correctness]\label{def:bs-correctness}
	A blind signature scheme is correct if there exists \(\negl\) such that for all messages \(\mu\),
	\[
	\Pr\!\left[
	\begin{array}{l}
		\PublicParamsBS \gets \Setup(\SecParam); \\
		(\vk,\sk)\gets \IKeyGen(\PublicParamsBS);\\
		\sig\gets \Holder^{\Issue: \langle \Issuer(\sk): \cdot  \rangle }(\mu,\vk);\\
		\Verify(\vk,\mu,\sig)=1
	\end{array}
	\right]
	\ge 1-\negl(\lambda).
	\]
\end{definition}


\begin{definition}[Blindness]\label{def:bs-blindness}
	A blind signature scheme is blind if there exists \(\negl\) such that for any PPT three-stage stateful adversary
	\(\adversary=(\adversary_1,\adversary_2,\adversary_3)\) and any two messages \(\mu_0,\mu_1\),
	\[
	\left|\Pr\!\left[
	\begin{array}{l}
		\PublicParamsBS\gets \Setup(\SecParam);\\
		\vk\gets \adversary_1(\PublicParamsBS);\\ b\getsunif\{0,1\};\\
		\adversary_2^{\Issue:\langle \cdot: \Holder(\mu_b,\vk)\rangle ,~ \Issue:\langle \cdot: \Holder(\mu_{1-b},\vk)\rangle }; \\
		\Holder(\mu_b,\vk) \text{ outputs }\sig_b;\\
		\Holder(\mu_{1-b},\vk) \text{ outputs }\sig_{1-b}\\
		\text{if }\bot\in\{\sig_0,\sig_1\}\text{ then }(\sig_0,\sig_1)\gets(\bot,\bot):\\
		b=\adversary_3(\sig_0,\sig_1)
	\end{array}
	\right]-\frac{1}{2}\right|
	\le \negl(\lambda).
	\]
\end{definition}



\begin{definition}[One-more unforgeability]\label{def:bs-euf-cma}
A blind signature scheme is one-more unforgeable if there exists \(\negl\) such that for any PPT adversary
\(\adversary^{\Issue: \langle \Issuer(\sk): \cdot \rangle}(\vk)\) that makes at most \(Q=Q(\lambda)\) oracle queries to \(\Issue\),
\[
  \Pr\!\left[
    \begin{array}{l}
     	\PublicParamsBS \gets \Setup(\SecParam); \\
      (\vk,\sk)\gets \IKeyGen(\PublicParamsBS);\\
      \{(\mu_i,\sig_i)\}_{i\in[Q+1]}\gets \adversary^{\Issue \langle \Issuer(\sk):\cdot \rangle}(\vk):\\
      \forall\,1\le i<j\le Q+1:\ \mu_i\neq \mu_j\ \land\ \forall\,i\in[Q+1]:\ \Verify(\vk,\mu_i,\sig_i)=1
    \end{array}
  \right]
  \le \negl(\lambda).
\]
\end{definition}

%Our issuance protocol (Section~\ref{sec:blind-signature}) instantiates this primitive and additionally packages the
%interaction with pre-sign and post-sign NIZKs tailored to ARC.

\subsection{Rate Limiting (with blind signatures)}
\lena{$\credential$ could also be $\sig$}
Rate limiting allows a holder $\Holder$ to obtain a credential $\credential$ that is valid under the verification key of an issuer $\Issuer$, which can be shown a limited number of times $\ell$.  
A blind signature algorithm allows a holder $\Holder$ to obtain signatures valid under the verification key of an issuer $\Issuer$, without the latter learning the signed message $\mu$. 
Our definitions for rate limiting follow \cite{C:Fischlin06}, with an additonal rate limit $\ell$ that limits the number of showings. 
When any algorithm $\mathcal{B}$ uses the interactive protocol $\Issue$ with $\Issuer$, we will write $\mathcal{B}^{\Issue:\langle \Issuer(\cdot):\cdot \rangle}$ and adjust this the notation for other interactions accordingly.


\begin{definition}[Rate Limiting]\label{def:rate-limit}
  A rate limting scheme is a tuple of algorithms 
 	\[
 	\RateLimit=(\Setup,\IKeyGen,\Issue,\Verify)
 	\]
 	such that:
 	\begin{itemize}[leftmargin=1.25em]
 		\item $\Setup(\SecParam)$ is a PPT algorithm that outputs public parameters $\PublicParamsRL$. All other algorithms obtain $\SecParam,\PublicParamsRL$ as implicit inputs.
 		\item $\IKeyGen()$ is a PPT algorithm that outputs issuer keys $(\vk,\sk)$.
  \item $\Issue$ is an interactive protocol between $\Issuer$ and
    $\Holder$, where $\Issuer$ has input $\sk$ and a rate limit $\ell$, and
    $\Holder$ has input $\mu,\vk$. At the end of the protocol, $\Holder$ obtains
    $\credential$ which can be shown at most $\ell$ times.
 		\item $\Verify(\vk,\mu,\credential,\nonce)$ is a deterministic polynomial-time algorithm that outputs a bit.
  The showings are unlinkable if they are correct. 
 	\end{itemize}
\end{definition}

\begin{definition}[Correctness]\label{def:rate-limit-correctness}
  %You can produce at least $\ell$ showings
	A rate limiting scheme is correct if there exists \(\negl\) such that for all messages \(\mu\), with a fresh \(\nonce\) within the rate limit $\ell$,
 	\[
 	\Pr\!\left[
 	\begin{array}{l}
 		\PublicParamsRL \gets \Setup(\SecParam); \\
 		(\vk,\sk)\gets \IKeyGen(\PublicParamsRL);\\
 		\sig\gets \Holder^{\Issue: \langle \Issuer(\sk,\ell): \cdot  \rangle }(\mu,\vk);\\
    \Verify(\vk,\mu,\sig,\nonce \in \{0,\ell\})=1
 	\end{array}
 	\right]
 	\ge 1-\negl(\lambda).
 	\]
\end{definition}

\lena{blindness would be privacy, idk if this is a good idea}
\begin{definition}[Privacy]\label{def:bs-blindness}
	A rate limiting scheme is private if there exists \(\negl\) such that for any PPT three-stage stateful adversary
	\(\adversary=(\adversary_1,\adversary_2,\adversary_3)\) and any two messages \(\mu_0,\mu_1\),
	\[
	\left|\Pr\!\left[
	\begin{array}{l}
		\PublicParamsBS\gets \Setup(\SecParam);\\
		\vk\gets \adversary_1(\PublicParamsBS);\\ b\getsunif\{0,1\};\\
		\adversary_2^{\Issue:\langle \cdot: \Holder(\mu_b,\vk)\rangle ,~ \Issue:\langle \cdot: \Holder(\mu_{1-b},\vk)\rangle }; \\
		\Holder(\mu_b,\vk) \text{ outputs }\sig_b;\\
		\Holder(\mu_{1-b},\vk) \text{ outputs }\sig_{1-b}\\
		\text{if }\bot\in\{\sig_0,\sig_1\}\text{ then }(\sig_0,\sig_1)\gets(\bot,\bot):\\
		b=\adversary_3(\sig_0,\sig_1)
	\end{array}
	\right]-\frac{1}{2}\right|
	\le \negl(\lambda).
	\]
\end{definition}


\lena{this is already nonce-defined, maybe we want to keep the definition more general. This is basically one-more unforgeability with the additional constraint that we can only spend  a credential once- should we make this two different defintions?}
\begin{definition}[Soundness]\label{def:rate-limit-soundness}
  % You can produce at most $\ell$ showings
	A rate limiting scheme is sound if there exists \(\negl\) such that for all messages \(\mu\), with a reused \(\nonce\) which is in the set of previously spent nonces \(\usednonce\) within the rate limit $\ell$ or a \(\nonce \geq \ell\),
 	\[
 	\Pr\!\left[
 	\begin{array}{l}
 		\PublicParamsRL \gets \Setup(\SecParam); \\
 		(\vk,\sk)\gets \IKeyGen(\PublicParamsRL);\\
 		\sig\gets \Holder^{\Issue: \langle \Issuer(\sk,\ell): \cdot  \rangle }(\mu,\vk);\\
    \Verify(\vk,\mu,\sig,(\nonce \geq \ell \lor \nonce \in \usednonce) )=0; \\
      \{(\mu_i,\sig_i)\}_{i\in[Q+1]}\gets \adversary^{\Issue \langle \Issuer(\sk):\cdot \rangle}(\vk):\\
      \forall\,1\le i<j\le Q+1:\ \mu_i\neq \mu_j\ \land\ \forall\,i\in[Q+1]:\ \Verify(\vk,\mu_i,\sig_i,\ell)=1 	\end{array}
 	\right]
 	\ge 1-\negl(\lambda).
 	\]
\end{definition}


\subsection{Anonymous Rate-Limiting Credentials (ARCs)}
\label{subsec:arc-security}
% LENA: I left the partial reveals out for now. 
Anonymous Rate-Limited credentials offer a similar syntax to Blind signatures\lena{/rate limiting}. 
However, they additionally provide an algorithm $\Show$ that may create multiple presentations of a previous output of $\Issue$.
We also explicitly introduce the additional verifier party $\Verifier$ that will obtain the presentations. Our defintion follows established syntax and security models~\cite{JOC:FucHanSla19,C:BLNS23} for anonymous credentials. 
We simplify the defintions, removing organizations and attributes which are not currently supported by ARC, and add a seperation between $\Setup$ and $\IKeyGen$.
The seperation is necessary to ensure the commitment matrix is chosen at random, which may lead to the commitment not being hiding. 
We also seperate $\Show$ and $\Verify$ to explicitly state that $\Show$ is the
client routine and $\Verify$ is the server routine of the credential
presentation. 

The blindness of ARC follows from AC unlinkability. 
At redemption, the same credential with two different nonces, and two different credentials with the same nonce, and two different credentials with different nonces, are indistinguishable to the verifier. 
\begin{definition}[ARC scheme syntax]\label{def:arc-syntax}
Let $\mathcal{N}$ be a finite set of nonces. 
An anonymous rate-limited credential (ARC) scheme is a tuple of algorithms
\[
  \ARC=(\Setup,\IKeyGen,\Issue,\Show,\Verify)
\]
such that: %\cb{should there also be a message attached to the credential?} \lena{could, but we don't have to add it in this revision imho}
\begin{itemize}[leftmargin=1.25em]
   \item $\Setup(\SecParam)$ is a PPT algorithm that outputs public parameters $\PublicParamsARC$. All other algorithms obtain $\SecParam,\PublicParamsARC$ as implicit inputs.

    \item $\IKeyGen()$ is a PPT algorithm that outputs issuer keys $(\vk,\sk)$.
    \item $\Issue(\cdot)$ is an interactive protocol between $\Issuer$ and $\Holder$, where $\Issuer$ has input $\sk$ and $\Holder$ has input $\vk$. At the end of the protocol, $\Holder$ obtains $\credential$.
  \item $\Show(\cdot)$ on input a credential $\credential$ and a nonce $\nonce\in \mathcal{N}$ outputs a presentation $\presentation$. 
  
  \item $\Verify(\vk,\presentation,\nonce)$ is a deterministic polynomial-time algorithm that outputs a bit.
 \end{itemize}
\end{definition}

Adjusting the correctness definition \ref{def:bs-correctness}, blindness definition \ref{def:bs-blindness} and unforgeability definition \ref{def:bs-euf-cma} to ARC follows trivially. However, we additionally require that if the same credential $\credential$ is presented multiple times, but for different nonces from $\mathcal{N}$, then the presentations should be unlinkable. To formalize this, we let a challenger $\Challenger$ create two credentials in interactions where the adversary controls the issuer. Then the challenger either creates two presentations for the same credential, or for different credentials.

\begin{definition}[Presentation unlinkability]\label{def:pres-unlink}
Let $\ARC$ be an ARC scheme. Consider the following experiment
$\mathsf{Exp}^{\mathsf{pres\mbox{-}unlink}}_{\ARC,\Adv}(\lambda)$ between a challenger $\Challenger$ and a two-stage stateful adversary $\Adv=(\Adv_1,\Adv_2)$:
\begin{enumerate}[leftmargin=1.25em]
  \item $\Challenger$ runs $\PublicParamsARC\gets \Setup(\SecParam)$, $(\vk,\sk)\gets \IKeyGen(\PublicParamsARC)$. 
\item Run $\adversary_1^{\Issue:\langle \cdot: \Challenger(\vk)\rangle ,~ \Issue:\langle \cdot: \Challenger(\vk)\rangle }$, at the end of which $\Challenger$ obtains $\credential_0,\credential_1$ and $\Adv_1$ outputs $\nonce_0,\nonce_1$.
\item If $\credential_0$ or $\credential_1$ is $\bot$ or if $\nonce_0=\nonce_1$ then output a random bit.

  \item $\Challenger$ samples $b\getsunif\{0,1\}$ and computes
  \[
    \presentation_0 \leftarrow \Show(\credential_0,\nonce_0),\qquad
    \presentation_1 \leftarrow \Show(\credential_b,\nonce_1).
  \]
  \item $\Adv_2(\presentation_0,\presentation_1)$ outputs a bit $b'$.
\end{enumerate}
Define $\Adv^{\mathsf{pres\mbox{-}unlink}}_{\ARC}(\Adv) := \big| \Pr[b'=b]-\tfrac12 \big|$.
We say ARC has \emph{presentation unlinkability} if for all PPT $\Adv$ this advantage is negligible in $\lambda$.

\end{definition}

%\paragraph{Nonce reuse caveat (intended).}
%If a nonce is reused, deterministic presentation components (e.g., a
%rate-limiting tag) breaks unlinkability for those presentations; therefore the experiment enforces nonce freshness.

\subsection{Pseudorandom functions}
\label{subsec:prelim-prf}
\label{subsec:prf}


\begin{definition}[Pseudorandom function (PRF)]\label{def:prf}
Let $\mathcal{K}$, $\mathcal{D}$ and $\mathcal{T}$ be sets such that $|\mathcal{K}|$ is exponential in $\lambda$.
A function family $\PRF:\mathcal{K}\times\mathcal{D}\to \mathcal{T}$ is a PRF if for all PPT adversaries
$\Adv$,
\[
  \Adv^{\mathsf{prf}}_{\PRF}(\Adv)
  := \Big|\Pr[\Adv^{\mathcal{O}_0}=1]-\Pr[\Adv^{\mathcal{O}_1}=1]\Big|
  \le \negl(\lambda),
\]
where $\mathcal{O}_0(\cdot)=\PRF(k,\cdot)$ for uniform $k\getsunif \mathcal{K}$ and $\mathcal{O}_1(\cdot)$ is a uniform
random function from $\mathcal{D}$ to $\mathcal{T}$.
\end{definition}

%\paragraph{ARC-SPRUCE instantiation rule.}
%In ARC-SPRUCE, the PRF key is the secret message component \(k\) inside the signed message \(m=\mu\Vert k\), and the
%PRF input includes the public nonce (and, if desired, a policy context encoding).

In this work, we instantiate PRFs based on the Poseidon2 construction~\cite{AFRICACRYPT:GraKhoSch23}.

It is based on the Poseidon2 permutation, which is defined as follows:
\begin{definition}[Poseidon2-style permutation structure]\label{def:poseidon-perm}
	Let $q$ be a prime, $d\ge 3$ such that $\gcd(d,q-1)=1$ and $t,R_F,R_P\in \mathbb{N}$ where $R_F$ is even.
	We call $R_F$ the \emph{external} rounds and $R_P$ the  \emph{internal} rounds. Furthermore, let
	 $\{c^{(i)}_j\}_{i\in [R_F]}^{j\in [t]}\subset\Fq$, $\{\hat c^{(i)}\}_{i\in [R_P]}\subset\Fq$ be public round constants and $\mM_E,\mM_I\in\Fq^{t\times t}$ be invertible maps.
	
	For $\vecx=(x_1,\dots,x_t)\in\Fq^t$, define external rounds
	\[
	E_i(\vecx) := \mM_E \cdot \big((x_1+c^{(i)}_1)^d,\dots,(x_t+c^{(i)}_t)^d\big)^\top,
	\]
	and internal rounds
	\[
	I_i(\vecx) := \mM_I \cdot \big((x_1+\hat c^{(i)})^d,\ x_2,\dots,x_t\big)^\top.
	\]
	Then
	\[
	P := E_{R_F}\circ\cdots\circ E_{R_F/2+1}\circ I_{R_P}\circ\cdots\circ I_{1}\circ E_{R_F/2}\circ\cdots\circ E_1.
	\]
\end{definition}

We additionally define
$\Tr_i:\Fq^t\to \Fq^{i}$ outputs the first $i$ elements of the input vector.

\begin{definition}[Poseidon2 PRF]\label{def:poseidon-prf}
Let $\ell_{\mathsf{tag}},\ell_{\mathsf{key}},\ell_{\mathsf{in}}\in \mathbb{N}$ and set $t=\ell_{\mathsf{key}}+\ell_{\mathsf{nonce}}$. For a key 
 $\veck\in\Fq^{\ell_{\mathsf{key}}}$  and public input 
$\vecn \in\Fq^{\ell_{\mathsf{in}}}$, we define the Poseidon2 PRF as 
\[
  \PRF(\veck,\vecn) := \Tr_{\ell_{\mathsf{tag}}}\big(P(\veck\Vert \vecn) + (\veck\Vert \vecn)\big),
\]
where $P:\Fq^t\to \Fq^t$ is the permutation from Definition \ref{def:poseidon-perm}.
\end{definition}



\paragraph{Constraint view (what the post-sign NIZK must prove).}
\cb{All of this should go to a later section.}
To prove $\mathsf{tag}=\PRF(k,\mathsf{nonce})$ inside SmallWood, where \(k\) is the signed secret part of
\(m=\mu\Vert k\), we arithmetize:
(i) S-box constraints of the form $z = y^d$ at selected checkpoints,
(ii) linear constraints for adding round constants and applying $M_E,M_I$ between checkpoints,
and (iii) feed-forward plus truncation constraints binding the public $\mathsf{tag}$ to the final state.
The constraint families are stated in Section~\ref{subsec:prf-checkpointed}, while checkpoint scheduling and extended
notes are in Appendix~\ref{app:prf}.
